% Auto-generated notebook appendix
\begin{tcolorbox}[colback=blue!5!white,colframe=blue!75!black,title=\textbf{Jupyter Notebook: Assignment\_1\_OLS.ipynb}]
โค้ดและผลลัพธ์ทั้งหมดด้านล่างเป็นการรันจริงจาก Jupyter Notebook
\end{tcolorbox}

\vspace{0.5em}
\noindent\textbf{Markdown:}
\begin{lstlisting}[style=output]
# CPE 342 Machine Learning — Assignment 1: Training Models

**67070501042 วิศิษฐ์ สุวรรณเนาว์**

OLS (Ordinary Least Squares) Linear Regression — Advertising Budget vs Sales
\end{lstlisting}
\vspace{0.5em}
\noindent\textbf{Markdown:}
\begin{lstlisting}[style=output]
## 0. Imports and Setup
\end{lstlisting}
\vspace{1em}
\noindent\textbf{In [1]:}
\begin{lstlisting}[style=pycode]
import numpy as np
import matplotlib
matplotlib.use('Agg')
import matplotlib.pyplot as plt
import matplotlib.font_manager as fm
import os
from scipy import stats

# Load Sarabun font for plots
_font_dir = os.path.join(os.environ.get('LOCALAPPDATA', ''), 'Microsoft', 'Windows', 'Fonts')
_sarabun_r = os.path.join(_font_dir, 'Sarabun-Regular.ttf')
_sarabun_b = os.path.join(_font_dir, 'Sarabun-Bold.ttf')

def _prop(bold=False, size=11):
    path = _sarabun_b if bold else _sarabun_r
    return fm.FontProperties(fname=path, size=size) if os.path.exists(path) else None

plt.style.use('seaborn-v0_8-whitegrid')
print("Libraries loaded.")
\end{lstlisting}
\smallskip
\noindent\textbf{Out[1]:}
\begin{lstlisting}[style=output]
Libraries loaded.
\end{lstlisting}
\vspace{0.5em}
\noindent\textbf{Markdown:}
\begin{lstlisting}[style=output]
## 1. Data
\end{lstlisting}
\vspace{1em}
\noindent\textbf{In [2]:}
\begin{lstlisting}[style=pycode]
# Advertising budget (X, in $1,000s) and Sales (Y, in 1,000 units)
X = np.array([3., 5., 2., 7., 8., 1., 4., 6., 9., 10.])
Y = np.array([6., 9., 4., 10., 12., 3., 7., 8., 13., 15.])
n = len(X)
months = np.arange(1, n + 1)

print(f"n = {n}")
print(f"X = {X}")
print(f"Y = {Y}")
\end{lstlisting}
\smallskip
\noindent\textbf{Out[2]:}
\begin{lstlisting}[style=output]
n = 10
X = [ 3.  5.  2.  7.  8.  1.  4.  6.  9. 10.]
Y = [ 6.  9.  4. 10. 12.  3.  7.  8. 13. 15.]
\end{lstlisting}
\vspace{0.5em}
\noindent\textbf{Markdown:}
\begin{lstlisting}[style=output]
## 2. Task 1: OLS Regression

### 2a. Setting up the Normal Equations

Minimise $S(\alpha, \beta) = \sum_i (Y_i - \alpha - \beta X_i)^2$ gives the Normal Equations:

$$\begin{bmatrix} n & \sum X_i \\ \sum X_i & \sum X_i^2 \end{bmatrix} \begin{bmatrix} \alpha \\ \beta \end{bmatrix} = \begin{bmatrix} \sum Y_i \\ \sum X_i Y_i \end{bmatrix}$$
\end{lstlisting}
\vspace{1em}
\noindent\textbf{In [3]:}
\begin{lstlisting}[style=pycode]
sum_X   = X.sum()
sum_Y   = Y.sum()
sum_XY  = (X * Y).sum()
sum_X2  = (X ** 2).sum()

print("Summation table:")
print(f"  sum(X)  = {sum_X}")
print(f"  sum(Y)  = {sum_Y}")
print(f"  sum(XY) = {sum_XY}")
print(f"  sum(X2) = {sum_X2}")
print()
print("Normal Equation (matrix form):")
print(f"  [{n:2d}  {sum_X:5.1f}] [alpha]   [{sum_Y:5.1f}]")
print(f"  [{sum_X:2.0f} {sum_X2:5.1f}] [beta ] = [{sum_XY:5.1f}]")
\end{lstlisting}
\smallskip
\noindent\textbf{Out[3]:}
\begin{lstlisting}[style=output]
Summation table:
  sum(X)  = 55.0
  sum(Y)  = 87.0
  sum(XY) = 583.0
  sum(X2) = 385.0

Normal Equation (matrix form):
  [10   55.0] [alpha]   [ 87.0]
  [55 385.0] [beta ] = [583.0]
\end{lstlisting}
\vspace{0.5em}
\noindent\textbf{Markdown:}
\begin{lstlisting}[style=output]
### 2b. Solving for Coefficients

**Method 0 — Closed-form via $S_{xx}$, $S_{xy}$** (derived directly from Normal Equations):

$$\beta = \frac{S_{xy}}{S_{xx}}, \qquad \alpha = \bar{Y} - \beta \bar{X}$$
\end{lstlisting}
\vspace{1em}
\noindent\textbf{In [4]:}
\begin{lstlisting}[style=pycode]
X_bar = X.mean()   # 5.5
Y_bar = Y.mean()   # 8.7

Sxx = np.sum((X - X_bar) ** 2)           # 82.5
Sxy = np.sum((X - X_bar) * (Y - Y_bar))  # 104.5

beta  = Sxy / Sxx                         # 1.266667
alpha = Y_bar - beta * X_bar              # 1.733333

print(f"X_bar = {X_bar},  Y_bar = {Y_bar}")
print(f"Sxx = {Sxx},  Sxy = {Sxy}")
print(f"beta  (slope)     = {beta:.6f}")
print(f"alpha (intercept) = {alpha:.6f}")
print(f"Fitted line: Y_hat = {alpha:.4f} + {beta:.4f} * X")
assert np.isclose(beta, 19/15)
assert np.isclose(alpha, 26/15)
\end{lstlisting}
\smallskip
\noindent\textbf{Out[4]:}
\begin{lstlisting}[style=output]
X_bar = 5.5,  Y_bar = 8.7
Sxx = 82.5,  Sxy = 104.5
beta  (slope)     = 1.266667
alpha (intercept) = 1.733333
Fitted line: Y_hat = 1.7333 + 1.2667 * X
\end{lstlisting}
\vspace{0.5em}
\noindent\textbf{Markdown:}
\begin{lstlisting}[style=output]
**Verification — Method 1: Cramer's Rule**
\end{lstlisting}
\vspace{1em}
\noindent\textbf{In [5]:}
\begin{lstlisting}[style=pycode]
A = np.array([[n, sum_X], [sum_X, sum_X2]], dtype=float)
b = np.array([sum_Y, sum_XY], dtype=float)

det_A       = np.linalg.det(A)
det_alpha   = np.linalg.det(np.column_stack([b,  A[:, 1]]))
det_beta    = np.linalg.det(np.column_stack([A[:, 0], b]))

alpha_cr = det_alpha / det_A
beta_cr  = det_beta  / det_A
print(f"Cramer's Rule → alpha = {alpha_cr:.6f},  beta = {beta_cr:.6f}")
assert np.isclose(alpha_cr, alpha) and np.isclose(beta_cr, beta)
print("✓ Matches Sxx/Sxy result")
\end{lstlisting}
\smallskip
\noindent\textbf{Out[5]:}
\begin{lstlisting}[style=output]
Cramer's Rule → alpha = 1.733333,  beta = 1.266667
✓ Matches Sxx/Sxy result
\end{lstlisting}
\vspace{0.5em}
\noindent\textbf{Markdown:}
\begin{lstlisting}[style=output]
**Verification — Method 2: Pseudoinverse $(A^T A)^{-1} A^T Y$**
\end{lstlisting}
\vspace{1em}
\noindent\textbf{In [6]:}
\begin{lstlisting}[style=pycode]
A_mat = np.column_stack([np.ones(n), X])   # design matrix [1 | X]
coeffs = np.linalg.inv(A_mat.T @ A_mat) @ A_mat.T @ Y
alpha_ps, beta_ps = coeffs
print(f"Pseudoinverse   → alpha = {alpha_ps:.6f},  beta = {beta_ps:.6f}")
assert np.isclose(alpha_ps, alpha) and np.isclose(beta_ps, beta)
print("✓ Matches all previous results")
\end{lstlisting}
\smallskip
\noindent\textbf{Out[6]:}
\begin{lstlisting}[style=output]
Pseudoinverse   → alpha = 1.733333,  beta = 1.266667
✓ Matches all previous results
\end{lstlisting}
\vspace{0.5em}
\noindent\textbf{Markdown:}
\begin{lstlisting}[style=output]
### 2c. Fitted Line
\end{lstlisting}
\vspace{1em}
\noindent\textbf{In [7]:}
\begin{lstlisting}[style=pycode]
Y_hat     = alpha + beta * X
residuals = Y - Y_hat

print(f"Fitted line: Y_hat = {alpha:.4f} + {beta:.4f} * X")
print()
print(f"{'Month':>5} {'X':>5} {'Y':>5} {'Y_hat':>8} {'e_i':>8}")
print("-" * 40)
for i in range(n):
    print(f"  {i+1:2d}   {X[i]:4.1f}  {Y[i]:4.1f}   {Y_hat[i]:7.4f}   {residuals[i]:+7.4f}")
\end{lstlisting}
\smallskip
\noindent\textbf{Out[7]:}
\begin{lstlisting}[style=output]
Fitted line: Y_hat = 1.7333 + 1.2667 * X

Month     X     Y    Y_hat      e_i
----------------------------------------
   1    3.0   6.0    5.5333   +0.4667
   2    5.0   9.0    8.0667   +0.9333
   3    2.0   4.0    4.2667   -0.2667
   4    7.0  10.0   10.6000   -0.6000
   5    8.0  12.0   11.8667   +0.1333
   6    1.0   3.0    3.0000   +0.0000
   7    4.0   7.0    6.8000   +0.2000
   8    6.0   8.0    9.3333   -1.3333
   9    9.0  13.0   13.1333   -0.1333
  10   10.0  15.0   14.4000   +0.6000
\end{lstlisting}
\vspace{0.5em}
\noindent\textbf{Markdown:}
\begin{lstlisting}[style=output]
## 3. Task 2: Interpretation

### 3a. Meaning of Slope and Intercept

- **Slope (β ≈ 1.2667):** For each additional \$1,000 in advertising budget, sales are predicted to increase by approximately **1,267 units**. This is an association in the sample, not proof of causation.
- **Intercept (α ≈ 1.7333):** At zero advertising budget (X = 0), predicted sales ≈ **1,733 units**. X = 0 is outside the observed range (1–10), so this is an extrapolated baseline.

### 3b. Prediction at X = 12 (\$12,000 budget)
\end{lstlisting}
\vspace{1em}
\noindent\textbf{In [8]:}
\begin{lstlisting}[style=pycode]
X_pred = 12.0
Y_pred = alpha + beta * X_pred

print(f"X = {X_pred} (i.e., $12,000 advertising budget)")
print(f"Y_hat = {alpha:.4f} + {beta:.4f} * {X_pred} = {Y_pred:.4f} thousand units")
print(f"      ≈ {Y_pred * 1000:,.0f} units")
print()
print("NOTE: X = 12 is beyond the observed maximum (X_max = 10).")
print("This prediction is extrapolation — treat with caution.")
assert np.isclose(Y_pred, 16.9333333333)
\end{lstlisting}
\smallskip
\noindent\textbf{Out[8]:}
\begin{lstlisting}[style=output]
X = 12.0 (i.e., $12,000 advertising budget)
Y_hat = 1.7333 + 1.2667 * 12.0 = 16.9333 thousand units
      ≈ 16,933 units

NOTE: X = 12 is beyond the observed maximum (X_max = 10).
This prediction is extrapolation — treat with caution.
\end{lstlisting}
\vspace{0.5em}
\noindent\textbf{Markdown:}
\begin{lstlisting}[style=output]
## 4. Task 3: Model Evaluation

### 4a. R² Calculation

$$R^2 = 1 - \frac{SS_{\text{res}}}{SS_{\text{tot}}}, \quad SS_{\text{res}} = \sum e_i^2, \quad SS_{\text{tot}} = \sum (Y_i - \bar{Y})^2$$
\end{lstlisting}
\vspace{1em}
\noindent\textbf{In [9]:}
\begin{lstlisting}[style=pycode]
SS_res = np.sum(residuals ** 2)
SS_tot = np.sum((Y - Y_bar) ** 2)
R2     = 1 - SS_res / SS_tot

print(f"SS_res = {SS_res:.6f}")
print(f"SS_tot = {SS_tot:.4f}")
print(f"R²     = {R2:.6f}  ({R2*100:.2f}%)")
print()
print(f"Residual mean = {residuals.mean():.6f}  (≈ 0 → no bias)")
print(f"Residual std  = {residuals.std(ddof=1):.6f}")

assert np.isclose(SS_res, 3.7333333333)
assert np.isclose(SS_tot, 136.1)
assert np.isclose(R2, 0.9725691893)
print("\n✓ All assertions passed")
\end{lstlisting}
\smallskip
\noindent\textbf{Out[9]:}
\begin{lstlisting}[style=output]
SS_res = 3.733333
SS_tot = 136.1000
R²     = 0.972569  (97.26%)

Residual mean = 0.000000  (≈ 0 → no bias)
Residual std  = 0.644061

✓ All assertions passed
\end{lstlisting}
\vspace{0.5em}
\noindent\textbf{Markdown:}
\begin{lstlisting}[style=output]
### 4b. Interpretation

97.26% of the observed variation in sales is explained by the fitted linear relationship with advertising budget.

**Important caveats:**
- High R² does **not** establish that advertising *causes* higher sales.
- R² alone does not guarantee the model will predict well on new data.
- Residual mean ≈ 0 confirms the model is unbiased over the sample.
\end{lstlisting}
\vspace{0.5em}
\noindent\textbf{Markdown:}
\begin{lstlisting}[style=output]
## 5. Visualisations
\end{lstlisting}
\vspace{0.5em}
\noindent\textbf{Markdown:}
\begin{lstlisting}[style=output]
### 5a. Scatter Plot with Regression Line
\end{lstlisting}
\vspace{1em}
\noindent\textbf{In [10]:}
\begin{lstlisting}[style=pycode]
fig, ax = plt.subplots(figsize=(8, 5.5))

# Data points
ax.scatter(X, Y, color='#2563EB', s=90, edgecolors='white',
           linewidths=1.5, zorder=5, label='Observed data')

# Regression line (extended to X=13)
X_line = np.linspace(0, 13, 300)
ax.plot(X_line, alpha + beta * X_line, color='#DC2626', linewidth=2.2,
        label=f'$\hat{{Y}} = {alpha:.4f} + {beta:.4f}X$', zorder=4)

# Prediction point
ax.scatter([X_pred], [Y_pred], color='#16A34A', s=140, marker='D',
           edgecolors='white', linewidths=1.5, zorder=6,
           label=f'Prediction at X=12 ($\hat{{Y}}$={Y_pred:.2f})')
ax.plot([X_pred, X_pred], [0, Y_pred], color='#16A34A',
        linestyle='--', linewidth=1, alpha=0.6, zorder=3)
ax.plot([0, X_pred], [Y_pred, Y_pred], color='#16A34A',
        linestyle='--', linewidth=1, alpha=0.6, zorder=3)

# Labels on each data point
for i in range(n):
    ax.annotate(f'M{i+1}', (X[i], Y[i]),
                textcoords='offset points', xytext=(7, 6), fontsize=8.5, color='#374151')

ax.set_xlabel('Advertising Budget (X) [$1,000s]')
ax.set_ylabel('Sales (Y) [1,000 units]')
ax.set_title('OLS Linear Regression: Advertising Budget vs Sales', fontsize=13)
ax.legend(loc='upper left', framealpha=0.9)
ax.set_xlim(-0.5, 13.5)
ax.set_ylim(0, 18)
plt.tight_layout()
fig.savefig('scatter_regression.pdf', dpi=300, bbox_inches='tight')
plt.show()
print("Saved scatter_regression.pdf")
\end{lstlisting}
\smallskip
\noindent\textbf{Out[10]:}
\begin{lstlisting}[style=output]
<>:10: SyntaxWarning: invalid escape sequence '\h'
<>:15: SyntaxWarning: invalid escape sequence '\h'
<>:10: SyntaxWarning: invalid escape sequence '\h'
<>:15: SyntaxWarning: invalid escape sequence '\h'
C:\Users\Admin\AppData\Local\Temp\ipykernel_26040\3933166873.py:10: SyntaxWarning: invalid escape sequence '\h'
  label=f'$\hat{{Y}} = {alpha:.4f} + {beta:.4f}X$', zorder=4)
C:\Users\Admin\AppData\Local\Temp\ipykernel_26040\3933166873.py:15: SyntaxWarning: invalid escape sequence '\h'
  label=f'Prediction at X=12 ($\hat{{Y}}$={Y_pred:.2f})')
Saved scatter_regression.pdf
C:\Users\Admin\AppData\Local\Temp\ipykernel_26040\3933166873.py:34: UserWarning: FigureCanvasAgg is non-interactive, and thus cannot be shown
  plt.show()
\end{lstlisting}
\vspace{0.5em}
\noindent\textbf{Markdown:}
\begin{lstlisting}[style=output]
### 5b. Residual Diagnostics Panel (2×2)
\end{lstlisting}
\vspace{1em}
\noindent\textbf{In [11]:}
\begin{lstlisting}[style=pycode]
fig, axes = plt.subplots(2, 2, figsize=(12, 9))

# -- (1) Residuals vs Fitted --
ax = axes[0, 0]
ax.scatter(Y_hat, residuals, color='#2563EB', s=90,
           edgecolors='white', linewidths=1.2, zorder=4)
ax.axhline(0, color='#DC2626', linestyle='--', linewidth=1.8, zorder=3)
for i in range(n):
    ax.annotate(f'M{i+1}', (Y_hat[i], residuals[i]),
                textcoords='offset points', xytext=(6, 5), fontsize=8.5, color='#374151')
ax.set_xlabel('Fitted values ($\hat{Y}$) [1,000 units]')
ax.set_ylabel('Residual ($e_i$)')
ax.set_title('Residuals vs Fitted Values')

# -- (2) Residual Distribution + Normal curve --
ax = axes[0, 1]
ax.hist(residuals, bins=6, color='#2563EB', edgecolor='white', alpha=0.82,
        density=True, zorder=4, label='Residuals')
x_range = np.linspace(residuals.min() - 0.5, residuals.max() + 0.5, 200)
ax.plot(x_range, stats.norm.pdf(x_range, 0, residuals.std(ddof=1)),
        color='#DC2626', linewidth=2.2, label='Normal PDF', zorder=5)
ax.axvline(0, color='#374151', linestyle='--', linewidth=1.2, zorder=3)
ax.set_xlabel('Residual')
ax.set_ylabel('Density')
ax.set_title('Residual Distribution')
ax.legend()

# -- (3) Normal Q-Q Plot --
ax = axes[1, 0]
(osm, osr), (slope_qq, intercept_qq, _) = stats.probplot(residuals, dist='norm')
ax.scatter(osm, osr, color='#7C3AED', s=90, edgecolors='white', linewidths=1.2, zorder=4)
xq = np.array([osm.min(), osm.max()])
ax.plot(xq, slope_qq * xq + intercept_qq,
        color='#DC2626', linewidth=2.2, zorder=3)
ax.set_xlabel('Theoretical Quantiles')
ax.set_ylabel('Sample Quantiles')
ax.set_title('Normal Q-Q Plot')

# -- (4) Scale-Location --
ax = axes[1, 1]
sqrt_abs = np.sqrt(np.abs(residuals))
ax.scatter(Y_hat, sqrt_abs, color='#059669', s=90,
           edgecolors='white', linewidths=1.2, zorder=4)
for i in range(n):
    ax.annotate(f'M{i+1}', (Y_hat[i], sqrt_abs[i]),
                textcoords='offset points', xytext=(6, 4), fontsize=8.5, color='#374151')
ax.set_xlabel('Fitted values ($\hat{Y}$) [1,000 units]')
ax.set_ylabel('$\sqrt{|e_i|}$')
ax.set_title('Scale-Location')

fig.suptitle('Residual Diagnostics', fontsize=15, fontweight='bold', y=1.01)
plt.tight_layout()
fig.savefig('diagnostics_panel.pdf', dpi=300, bbox_inches='tight')
plt.show()
print("Saved diagnostics_panel.pdf")
\end{lstlisting}
\smallskip
\noindent\textbf{Out[11]:}
\begin{lstlisting}[style=output]
<>:11: SyntaxWarning: invalid escape sequence '\h'
<>:47: SyntaxWarning: invalid escape sequence '\h'
<>:48: SyntaxWarning: invalid escape sequence '\s'
<>:11: SyntaxWarning: invalid escape sequence '\h'
<>:47: SyntaxWarning: invalid escape sequence '\h'
<>:48: SyntaxWarning: invalid escape sequence '\s'
C:\Users\Admin\AppData\Local\Temp\ipykernel_26040\2831409773.py:11: SyntaxWarning: invalid escape sequence '\h'
  ax.set_xlabel('Fitted values ($\hat{Y}$) [1,000 units]')
C:\Users\Admin\AppData\Local\Temp\ipykernel_26040\2831409773.py:47: SyntaxWarning: invalid escape sequence '\h'
  ax.set_xlabel('Fitted values ($\hat{Y}$) [1,000 units]')
C:\Users\Admin\AppData\Local\Temp\ipykernel_26040\2831409773.py:48: SyntaxWarning: invalid escape sequence '\s'
  ax.set_ylabel('$\sqrt{|e_i|}$')
Saved diagnostics_panel.pdf
C:\Users\Admin\AppData\Local\Temp\ipykernel_26040\2831409773.py:54: UserWarning: FigureCanvasAgg is non-interactive, and thus cannot be shown
  plt.show()
\end{lstlisting}
\vspace{0.5em}
\noindent\textbf{Markdown:}
\begin{lstlisting}[style=output]
### 5c. Actual vs Fitted Sales
\end{lstlisting}
\vspace{1em}
\noindent\textbf{In [12]:}
\begin{lstlisting}[style=pycode]
fig, ax = plt.subplots(figsize=(6.5, 6))

sc = ax.scatter(Y, Y_hat, c=months, cmap='viridis', s=120,
               edgecolors='white', linewidths=1.2, zorder=4)
cbar = fig.colorbar(sc, ax=ax, label='Month')

lims = [min(Y.min(), Y_hat.min()) - 0.4, max(Y.max(), Y_hat.max()) + 0.4]
ax.plot(lims, lims, color='#DC2626', linestyle='--',
        linewidth=1.8, label='Perfect fit', zorder=3)

for i in range(n):
    ax.annotate(f'M{i+1}', (Y[i], Y_hat[i]),
                textcoords='offset points', xytext=(6, 4), fontsize=8.5, color='#374151')

ax.set_xlabel('Actual Sales $Y$ [1,000 units]')
ax.set_ylabel('Fitted Sales $\hat{Y}$ [1,000 units]')
ax.set_title('Actual vs Fitted Sales', fontsize=13)
ax.legend(loc='upper left')
ax.set_xlim(lims); ax.set_ylim(lims)
ax.set_aspect('equal')
plt.tight_layout()
fig.savefig('actual_vs_fitted.pdf', dpi=300, bbox_inches='tight')
plt.show()
print("Saved actual_vs_fitted.pdf")
\end{lstlisting}
\smallskip
\noindent\textbf{Out[12]:}
\begin{lstlisting}[style=output]
<>:16: SyntaxWarning: invalid escape sequence '\h'
<>:16: SyntaxWarning: invalid escape sequence '\h'
C:\Users\Admin\AppData\Local\Temp\ipykernel_26040\2730993131.py:16: SyntaxWarning: invalid escape sequence '\h'
  ax.set_ylabel('Fitted Sales $\hat{Y}$ [1,000 units]')
Saved actual_vs_fitted.pdf
C:\Users\Admin\AppData\Local\Temp\ipykernel_26040\2730993131.py:23: UserWarning: FigureCanvasAgg is non-interactive, and thus cannot be shown
  plt.show()
\end{lstlisting}
\vspace{0.5em}
\noindent\textbf{Markdown:}
\begin{lstlisting}[style=output]
## 6. Summary of Results

| Task | Answer |
|------|--------|
| **1a** | Normal Equations: $10\alpha + 55\beta = 87$, $55\alpha + 385\beta = 583$ |
| **1b** | Intercept $\alpha \approx 1.7333$, Slope $\beta \approx 1.2667$ |
| **1c** | Fitted line: $\hat{Y} = 1.7333 + 1.2667X$ |
| **2a** | Slope: +\$1,000 ads → +1,267 units sold (association, not causation) |
| **2b** | At \$12,000 budget: $\hat{Y} \approx 16.93$ thousand units ≈ 16,933 units (extrapolation) |
| **3a** | $R^2 \approx 0.9726$ (97.26%) |
| **3b** | 97.26% of sales variation explained by advertising budget; residual mean = 0 (no bias) |
\end{lstlisting}